\section{Introducción y conceptos generales}
\label{sec:introduccion-conceptos}

La economía es la ciencia que estudia cómo las personas, las empresas, los hogares y los países toman decisiones cuando tienen recursos limitados y necesidades múltiples. La escasez es el punto de partida de la disciplina: no hay suficientes recursos para satisfacer todos los deseos al mismo tiempo. Por esa razón se debe elegir, y cada elección implica sacrificar otras alternativas.

\begin{concepto}{Escasez}
La escasez es la condición fundamental del mundo: los deseos y necesidades son prácticamente ilimitados, pero los recursos (tiempo, dinero, tierra, capital, trabajo) son limitados.
\end{concepto}

En Economía, esto produce una idea central: \textbf{no hay decisiones sin costo}. Incluso cuando no pagas dinero, estás renunciando a otra opción. Ese valor de la mejor alternativa sacrificada se llama \textbf{costo de oportunidad}.

\subsection{¿Qué es la economía?}

La palabra economía proviene del griego \textit{oikos} (casa) y \textit{nomos} (ley, administración). En su origen, estaba ligada a la administración del hogar. Con el tiempo, la economía pasó a estudiar también la administración de la riqueza, los recursos y las decisiones sociales.

\begin{ejemplo}
Si una persona tiene 10 horas disponibles, puede estudiar o trabajar. Si decide estudiar, el costo de oportunidad es el salario que dejó de ganar al no trabajar.
\end{ejemplo}

\subsection{Motivación para estudiar economía}

Estudiar economía permite responder preguntas fundamentales:
\begin{itemize}
    \item ¿Por qué hay pobreza y riqueza?
    \item ¿Por qué algunas políticas públicas ayudan y otras empeoran la situación?
    \item ¿Por qué los precios suben o bajan?
    \item ¿Qué determina el empleo, la inflación y el crecimiento?
\end{itemize}

La economía también ayuda a entender la conducta humana en términos de incentivos, decisiones y consecuencias.

\subsection{Microeconomía y macroeconomía}

\begin{itemize}
    \item \textbf{Microeconomía:} analiza el comportamiento de individuos, familias y empresas, y cómo interactúan en mercados específicos.
    \item \textbf{Macroeconomía:} estudia la economía en su conjunto: inflación, desempleo, crecimiento, dinero, deuda pública y política económica.
\end{itemize}

\begin{observacion}{Importante}
Aunque son ramas distintas, ambas están conectadas. Las decisiones individuales afectan agregados macroeconómicos, y los fenómenos macroeconómicos se originan en decisiones microeconómicas.
\end{observacion}

\subsection{Escasez, eficiencia, equidad y disyuntiva}

La escasez obliga a elegir. Esta elección siempre tiene un trade-off, es decir, un intercambio. La sociedad debe decidir cómo asignar recursos entre objetivos distintos.

\begin{itemize}
    \item \textbf{Eficiencia:} sacar el máximo provecho posible de los recursos escasos.
    \item \textbf{Eficacia:} lograr el objetivo deseado, aunque no se utilicen los recursos del modo más eficiente.
    \item \textbf{Equidad:} repartir de manera justa los beneficios o los costos.
\end{itemize}

\begin{ejemplo}
Un gobierno puede decidir aumentar el gasto en salud, pero eso significa menos dinero disponible para educación o infraestructura.
\end{ejemplo}

\subsection{Incentivos}

Los incentivos son cambios en costos o beneficios que alteran el comportamiento de las personas.

\begin{ejemplo}
Si el impuesto al alcohol aumenta, algunas personas dejan de consumirlo. Si un subsidio facilita la compra de bicicletas, más gente puede optar por ese transporte.
\end{ejemplo}

\subsection{Racionalidad}

En economía, se asume que los individuos actúan de manera racional, es decir, toman decisiones coherentes con sus objetivos, dadas las restricciones de información y recursos.

\begin{concepto}{Racionalidad}
La racionalidad significa que las personas eligen lo que consideran mejor para ellas, dadas las alternativas disponibles.
\end{concepto}

\subsection{Costo de oportunidad}

\begin{concepto}{Costo de oportunidad}
Es el valor de la mejor alternativa que se sacrifica al elegir otra opción.
\end{concepto}

El costo de oportunidad no siempre es monetario. Puede incluir tiempo, esfuerzo, oportunidades perdidas y otros beneficios.

\begin{ejemplo}
Si una persona estudia en vez de trabajar, el costo de oportunidad puede ser el salario que deja de ganar.
\end{ejemplo}

\begin{itemize}
    \item \textbf{Costo explícito:} dinero que se desembolsa directamente.
    \item \textbf{Costo implícito:} valor del sacrificio no monetario.
\end{itemize}

\subsection{Análisis marginal}

Los agentes no toman decisiones de manera absoluta, sino en términos de cambios pequeños, conocidos como \'márgenes\'.

\begin{concepto}{Beneficio marginal}
Es el beneficio extra que obtiene una persona al consumir o producir una unidad adicional.
\end{concepto}

\begin{concepto}{Costo marginal}
Es el costo extra de producir o consumir una unidad adicional.
\end{concepto}

\textbf{Regla marginalista:} se realiza una acción si el beneficio marginal es mayor o igual que el costo marginal.

\begin{ejemplo}
Una persona decide si estudiar una hora más: el beneficio marginal puede ser mayor si esa hora aumenta mucho la probabilidad de aprobar; si no, quizá no vale la pena.
\end{ejemplo}

\subsection{Fallas de mercado}

Aunque el mercado puede asignar eficientemente recursos, en algunas situaciones fallan sus mecanismos. Eso ocurre por ejemplos como:
\begin{itemize}
    \item poder de mercado (monopolio),
    \item externalidades,
    \item información asimétrica,
    \item bienes públicos y recursos comunes.
\end{itemize}

\begin{observacion}{Importante}
Un mercado puede ser eficiente en términos de asignación, pero incompleto o injusto en términos distributivos. Eso es un problema central de la economía pública.
\end{observacion}

\subsection{¿Qué es un modelo en economía?}

Un modelo es una simplificación de la realidad usada para estudiar un fenómeno de forma clara. Los economistas asumen condiciones hipotéticas para poder aislar ciertas relaciones.

\begin{ejemplo}
Para analizar la demanda de un bien, se asume que todo lo demás permanece constante, salvo el precio del bien en cuestión.
\end{ejemplo}

\subsection{Ceteris paribus}

\begin{concepto}{Ceteris paribus}
Expresión latina que significa “manteniendo todo lo demás constante”.
\end{concepto}

Este supuesto permite analizar un efecto aislado. Por ejemplo, si sube el precio de un bien, vemos cómo responde la cantidad demandada, manteniendo constantes los ingresos, gustos y los precios de otros bienes.

\subsection{Causalidad y correlación}

\begin{itemize}
    \item \textbf{Correlación:} dos variables se mueven juntas.
    \item \textbf{Causalidad:} una variable provoca un cambio en otra.
\end{itemize}

\begin{advertencia}{Error clásico}
La correlación no implica causalidad. Muchas veces hay una variable omitida que explica la relación.
\end{advertencia}

\begin{ejemplo}
Hay mayor consumo de helados cuando hace calor y también más personas se ahogan. No significa que los helados causen ahogamiento; ambos se relacionan con el clima.
\end{ejemplo}

\subsection{Economía positiva y normativa}

\begin{itemize}
    \item \textbf{Economía positiva:} describe el mundo tal como es. Busca explicar relaciones observables.
    \item \textbf{Economía normativa:} expresa juicios de valor sobre cómo debería ser el mundo.
\end{itemize}

\begin{ejemplo}
- Positiva: “Si aumenta el precio de la gasolina, la cantidad demandada cae.”
- Normativa: “El gobierno debería subsidiar el transporte para reducir la contaminación.”
\end{ejemplo}

\subsection{Reseña final del bloque de introducción}

La economía es el estudio de la elección bajo escasez. Sus herramientas básicas son: escasez, incentivos, costo de oportunidad, análisis marginal, modelos y supuestos. Todo esto permite entender cómo se deciden los recursos y cómo funcionan los mercados y la economía en general.

\begin{resumen}{Lo más importante del bloque}
- Hay escasez.
- Hay elecciones.
- Todo tiene costo de oportunidad.
- Los incentivos cambian el comportamiento.
- La economía usa modelos para estudiar la realidad.
- El objetivo es explicar y evaluar decisiones y resultados.
\end{resumen}
